
Data attribution methods aim to quantify the influence of individual training examples on model behavior. Related work is organized into three themes: foundational methods, scalable approximations, and evaluation protocols.

\subsubsection{2.1 Influence Functions and Their Limitations}

The modern era of data attribution began with Koh and Liang (2017), who adapted classical influence functions from robust statistics to deep learning. Their approach approximates the effect of removing a training point by computing the inverse Hessian-vector product, enabling attribution without expensive retraining. This foundational work demonstrated practical applications including debugging mislabeled examples and understanding model predictions.

Basu et al. (2020) showed that influence functions are fragile in deep networks---their accuracy varies with network depth, width, and regularization strength. This fragility was connected to the non-convexity of deep learning loss landscapes, where the Hessian approximation quality degrades in ways that are difficult to predict. While this work illuminated when influence functions fail, it did not characterize the trade-offs between different quality dimensions when methods produce reasonable but disagreeing results.

\subsubsection{2.2 Scalable Attribution Approximations}

The computational cost of exact influence functions motivated efficient approximations. TracIn \citep{pruthi2020tracin} bypasses Hessian computation by using gradient similarity across training checkpoints. TRAK \citep{park2023trak} employs random projections to reduce dimensionality while preserving attribution signal, reporting 0.7-0.9 rank correlation with LOO retraining at orders of magnitude speedup.

More recent work has targeted specific model architectures. DataInf \citep{kwon2023datainf} exploits the low-rank structure of LoRA fine-tuning to derive closed-form influence expressions. LoRIF \citep{li2026lorif} addresses I/O bottlenecks through low-rank factorization. MAGIC \citep{ilyas2025magic} combines classical methods with metadifferentiation.

Each of these methods reports improvements on specific metrics and datasets, but direct comparison is complicated by differing evaluation protocols, compute budgets, and quality metrics. The current work addresses this by normalizing computation via gradient-equivalent operations and measuring multiple quality dimensions simultaneously.

\subsubsection{2.3 Evaluation Protocols and Benchmarks}

Attribution evaluation typically follows one of two paradigms: correlation with ground truth or downstream task performance. Ground truth approaches compute expensive LOO retraining and measure rank correlation. Downstream approaches evaluate whether removing top-attributed examples degrades model performance.

Nguyen et al. (2023) raised concerns about evaluation reliability, showing that attribution estimates can be dominated by noise from random initialization and SGD stochasticity rather than genuine training data signal. This work identifies scenarios where attribution is unreliable but does not characterize the multi-objective structure when attribution is reliable.

A limitation of existing evaluation is single-metric focus. Papers typically report one headline number without characterizing the full quality profile. This obscures potential trade-offs.

\subsubsection{2.4 Positioning}

Prior work has advanced attribution efficiency, identified failure modes, and raised evaluation concerns. However, no work has systematically characterized the multi-objective structure of attribution quality or connected observed trade-offs to the underlying optimization geometry. This work provides a framework for understanding why methods disagree and how practitioners might select among them based on application requirements.
